\begin{definition}[Finite ideal-intersection predicate]
\label{def:finite-ideal-intersection-predicate}
Let $R$ carry a $\RingUp$ certificate with carrier $C$, classifier $\sim_R$, addition $+_R$, additive inverse $-_R$, zero $0_R$, and multiplication $\cdot_R$. Let $xs:\mathsf{HistSpine}$ be a finite BEDC-history spine, and let $\mathcal{I}$ assign to every position $i:\mathsf{Pos}(xs)$ a predicate $\mathcal{I}_i:\mathsf{BHist}\to\mathsf{Prop}$ over the same ring data. The finite relative intersection predicate is
$$
\mathsf{FinIdealMeet}_{C,xs,\mathcal{I}}(x)
\Longleftrightarrow
C(x)\land \forall i:\mathsf{Pos}(xs).\;\mathcal{I}_i(x).
$$
The carrier conjunct makes the empty finite spine denote the ambient carried ring predicate, while a nonempty spine is the usual pointwise intersection of the displayed ideal predicates.
\end{definition}
\leandef{BEDC.Derived.IdealUp.FiniteIdealMeet}

\begin{theorem}[Finite ideal intersection closure]
\label{thm:finite-ideal-intersection-closure}
Let $R$ carry a $\RingUp$ certificate as in \autoref{ch:concrete-instances-ring-namecert}, with carrier $C$, classifier $\sim_R$, addition $+_R$, additive inverse $-_R$, zero $0_R$, and multiplication $\cdot_R$. Let $xs:\mathsf{HistSpine}$ be a finite BEDC-history spine, and let $\mathcal{I}$ assign to every position $i:\mathsf{Pos}(xs)$ a predicate $\mathcal{I}_i:\mathsf{BHist}\to\mathsf{Prop}$. Assume that, for every $i:\mathsf{Pos}(xs)$, the predicate $\mathcal{I}_i$ is an $\IdealUp$ predicate over the same $\RingUp$ data, with the rows specified in \autoref{sec:ideal-certificate}: carrier support, zero membership, additive closure, additive inverse closure, internal multiplication closure, classifier transport, and left and right absorption by every carried ring endpoint. For
$$
K(x)\Longleftrightarrow \mathsf{FinIdealMeet}_{C,xs,\mathcal{I}}(x),
$$
the predicate $K$ is an $\IdealUp$ predicate over $R$. Explicitly, $K$ has carrier support and the rows
$$
K(0_R),\qquad
K(x)\land K(y)\Longrightarrow K(x+_R y),\qquad
K(x)\Longrightarrow K(-_R x),
$$
$$
K(x)\land K(y)\Longrightarrow K(x\cdot_R y),\qquad
K(x)\land x\sim_R y\Longrightarrow K(y),
$$
and, for every $r$ with $C(r)$,
$$
K(x)\Longrightarrow K(r\cdot_R x)\land K(x\cdot_R r).
$$
\end{theorem}
\leanchecked{BEDC.Derived.IdealUp.FiniteIdealMeet\_closure\_rows}
\begin{proof}
Unfold \autoref{def:finite-ideal-intersection-predicate}. The carrier support row is the first conjunct of $K(x)$.

For zero membership, the ring certificate supplies $C(0_R)$. Fix a position $i:\mathsf{Pos}(xs)$. The zero row in the $i$th ideal certificate gives $\mathcal{I}_i(0_R)$. Since $i$ was arbitrary, $0_R$ satisfies every predicate in the finite family, and together with $C(0_R)$ this is $K(0_R)$.

For additive closure, assume $K(x)$ and $K(y)$. The carrier conjuncts give $C(x)$ and $C(y)$, so the additive closure row of the ring certificate gives $C(x+_R y)$. For an arbitrary position $i$, the universal conjuncts of $K(x)$ and $K(y)$ give $\mathcal{I}_i(x)$ and $\mathcal{I}_i(y)$. The additive closure row of the $i$th ideal certificate gives $\mathcal{I}_i(x+_R y)$. Thus every indexed predicate holds at $x+_R y$, and the carrier witness just obtained repacks as $K(x+_R y)$.

For additive inverse closure, assume $K(x)$. The ring inverse-closure row gives $C(-_R x)$ from $C(x)$. For each position $i$, the inverse row of the $i$th ideal certificate gives $\mathcal{I}_i(-_R x)$ from $\mathcal{I}_i(x)$. Hence all finite-family predicates hold at $-_R x$, and $K(-_R x)$ follows.

For internal multiplication closure, assume $K(x)$ and $K(y)$. The ring multiplication-closure row gives $C(x\cdot_R y)$ from $C(x)$ and $C(y)$. For every position $i$, the universal conjuncts give $\mathcal{I}_i(x)$ and $\mathcal{I}_i(y)$, and the internal multiplication row of the $i$th ideal certificate gives $\mathcal{I}_i(x\cdot_R y)$. Therefore $K(x\cdot_R y)$.

For classifier transport, assume $K(x)$ and $x\sim_R y$. Carrier transport from the naming certificate underlying the $\RingUp$ data carries $C(x)$ along $x\sim_R y$ to obtain $C(y)$. For each position $i$, the classifier-transport row of the $i$th ideal certificate gives $\mathcal{I}_i(y)$ from $\mathcal{I}_i(x)$ and the same classifier witness. Pairing these facts gives $K(y)$.

For absorption, assume $K(x)$ and $C(r)$. The ring multiplication-closure row gives both $C(r\cdot_R x)$ and $C(x\cdot_R r)$. For each position $i$, left absorption in the $i$th ideal certificate gives $\mathcal{I}_i(r\cdot_R x)$ from $C(r)$ and $\mathcal{I}_i(x)$, and right absorption gives $\mathcal{I}_i(x\cdot_R r)$ from the same data. Hence the finite family holds at both absorbed endpoints. Together with the carrier witnesses, this gives $K(r\cdot_R x)$ and $K(x\cdot_R r)$.

These are exactly the carrier, additive sub-rng, multiplicative sub-rng, classifier-stability, and two-sided ambient absorption rows required of an $\IdealUp$ predicate over the fixed $\RingUp$ certificate.
\end{proof}

\begin{theorem}[Finite ideal meet quotient-ring kernel scope]
\label{thm:finite-ideal-meet-quotientring-kernel-scope}
Let $R$ carry a $\RingUp$ certificate with carrier $C$, classifier $\sim_R$, addition $+_R$, additive inverse $-_R$, zero $0_R$, and multiplication $\cdot_R$. Let $xs:\mathsf{HistSpine}$ be a finite BEDC-history spine, and let $\mathcal{I}$ assign to every position $i:\mathsf{Pos}(xs)$ a predicate $\mathcal{I}_i:\mathsf{BHist}\to\mathsf{Prop}$ over the same ring data. Assume that every indexed predicate $\mathcal{I}_i$ is an $\IdealUp$ predicate over these ring rows, and put
$$
K(x)\Longleftrightarrow \mathsf{FinIdealMeet}_{C,xs,\mathcal{I}}(x).
$$
Let $Q_K$ be the quotient-kernel source package of \autoref{def:ideal-quotient-kernel-source-package}. The package $Q_K$ has the same diagonal, endpoint-transport, and multiplication-descent scope as the quotient-kernel packages built from any other $\IdealUp$ predicate. Thus $\QuotientRingUp$ may consume finite intersections of ideals through the quotient-free kernel interface, with no finite-meet-specific quotient carrier or quotient primitive.
\end{theorem}
\begin{proof}
By \autoref{thm:finite-ideal-intersection-closure}, the predicate $K$ is an $\IdealUp$ predicate over the same $\RingUp$ certificate as the finite family $\mathcal{I}$. This supplies the carrier support, sub-rng, classifier-stability, and two-sided ambient absorption rows required by the quotient-kernel source package.

Apply \autoref{thm:ideal-quotient-kernel-diagonal-exactness} to $K$ to obtain diagonal exactness. Apply \autoref{thm:ideal-quotient-kernel-transport-surface} to $K$ to obtain endpoint transport along $\sim_R$ for carried endpoints.

Finally apply \autoref{thm:ideal-quotient-kernel-multiplication-descent} to $K$. The multiplication row is therefore the same product-difference membership row used by arbitrary $\IdealUp$ predicates, so the quotient-ring consumer reads the finite meet only through $Q_K(x,y)\Longleftrightarrow C(x)\land C(y)\land K(x-_R y)$.
\end{proof}

\begin{definition}[Ring map ideal preimage predicate]
\label{def:ring-map-ideal-preimage-predicate}
Let $S$ and $T$ carry $\RingUp$ certificates, with carriers $C_S,C_T$, classifiers $\sim_S,\sim_T$, zeros $0_S,0_T$, additions $+_S,+_T$, additive inverses $-_S,-_T$, and multiplications $\cdot_S,\cdot_T$. Let $J$ be an $\IdealUp$ carried predicate over $T$ in the sense of \autoref{sec:ideal-carrier}. For a function $f:S\to T$, the source-carried ideal preimage predicate is
$$
\mathsf{Preim}^{\mathrm{ideal}}_{f,J}(x)\Longleftrightarrow C_S(x)\land J(f(x)).
$$
The first component records source-carrier evidence, and the second component records membership of the image endpoint in the target ideal.
\end{definition}

\begin{theorem}[Ring map ideal preimage IdealUp closure]
\label{thm:ring-map-ideal-preimage-ideal-closure}
Let $S$ and $T$ carry $\RingUp$ certificates as in \autoref{ch:concrete-instances-ring-namecert}, with notation as in \autoref{def:ring-map-ideal-preimage-predicate}. Let $f:S\to T$ satisfy carrier preservation and classifier preservation on carried source endpoints:
$$
C_S(x)\Longrightarrow C_T(f(x))
$$
and
$$
C_S(x)\land C_S(y)\land x\sim_S y\Longrightarrow f(x)\sim_T f(y).
$$
Assume also that $f$ preserves zero, addition, additive inverse, and multiplication on carried source endpoints:
$$
f(0_S)\sim_T0_T,
$$
$$
C_S(x)\land C_S(y)\Longrightarrow f(x+_S y)\sim_T f(x)+_T f(y),
$$
$$
C_S(x)\Longrightarrow f(-_S x)\sim_T -_T f(x),
$$
and
$$
C_S(x)\land C_S(y)\Longrightarrow f(x\cdot_S y)\sim_T f(x)\cdot_T f(y).
$$
Let $J$ be an $\IdealUp$ predicate over $T$. Assume its ideal certificate supplies the rows specified in \autoref{sec:ideal-certificate}: carrier support, zero membership, additive closure, additive inverse closure, internal multiplication closure, classifier transport, and left and right absorption by every carried target ring endpoint. For $K:=\mathsf{Preim}^{\mathrm{ideal}}_{f,J}$ from \autoref{def:ring-map-ideal-preimage-predicate}, $K$ is an $\IdealUp$ predicate over the source ring $S$. Explicitly, $K$ has carrier support and the rows
$$
K(0_S),\qquad K(x)\land K(y)\Longrightarrow K(x+_S y),\qquad K(x)\Longrightarrow K(-_S x),
$$
$$
K(x)\land K(y)\Longrightarrow K(x\cdot_S y),\qquad K(x)\land x\sim_S y\Longrightarrow K(y),
$$
and, for every $r$ with $C_S(r)$,
$$
K(x)\Longrightarrow K(r\cdot_S x)\land K(x\cdot_S r).
$$
\end{theorem}
\leanchecked{BEDC.Derived.IdealUp.RingMapIdealPreimage\_closure\_rows}
\begin{proof}
Unfold \autoref{def:ring-map-ideal-preimage-predicate}. Carrier support is the first component of a $K$ witness.

For zero membership, the source $\RingUp$ certificate carries $0_S$. The target ideal zero row gives $J(0_T)$. The zero-preservation row gives $f(0_S)\sim_T0_T$, and target classifier symmetry from \autoref{def:lean-carrier-naming-certificate} gives $0_T\sim_T f(0_S)$. Target ideal classifier transport therefore gives $J(f(0_S))$, so $K(0_S)$.

For additive closure, assume $K(x)$ and $K(y)$. The source carrier components give $C_S(x)$ and $C_S(y)$, hence the source ring certificate carries $x+_S y$. The target ideal additive row gives $J(f(x)+_T f(y))$ from $J(f(x))$ and $J(f(y))$. The addition-preservation row gives $f(x+_S y)\sim_T f(x)+_T f(y)$, and target classifier symmetry reverses this comparison. Applying target ideal classifier transport gives $J(f(x+_S y))$. Together with the source carrier witness this is $K(x+_S y)$.

For additive inverse closure, assume $K(x)$. Source inverse closure carries $-_S x$. The target ideal inverse row gives $J(-_T f(x))$. The inverse-preservation row gives $f(-_S x)\sim_T -_T f(x)$, whose symmetric form transports $J(-_T f(x))$ to $J(f(-_S x))$. Thus $K(-_S x)$.

For internal multiplication closure, assume $K(x)$ and $K(y)$. Source multiplication closure carries $x\cdot_S y$. The target ideal internal multiplication row gives $J(f(x)\cdot_T f(y))$. The multiplication-preservation row gives $f(x\cdot_S y)\sim_T f(x)\cdot_T f(y)$, and target classifier symmetry orients this comparison for ideal transport. Hence $J(f(x\cdot_S y))$, so $K(x\cdot_S y)$.

For source classifier transport, assume $K(x)$ and $x\sim_S y$. Source carrier transport from \autoref{def:lean-carrier-naming-certificate} carries $C_S(x)$ along $x\sim_S y$ to obtain $C_S(y)$. Classifier preservation gives $f(x)\sim_T f(y)$. Target ideal classifier transport sends $J(f(x))$ along this comparison to $J(f(y))$. Therefore $K(y)$.

For left absorption, assume $K(x)$ and $C_S(r)$. Source multiplication closure carries $r\cdot_S x$, and carrier preservation gives $C_T(f(r))$. The target ideal left-absorption row gives $J(f(r)\cdot_T f(x))$ from $C_T(f(r))$ and $J(f(x))$. Multiplication preservation gives $f(r\cdot_S x)\sim_T f(r)\cdot_T f(x)$; after target classifier symmetry, ideal transport yields $J(f(r\cdot_S x))$. Thus $K(r\cdot_S x)$.

For right absorption, the same data carry $x\cdot_S r$ in the source, and the target ideal right-absorption row gives $J(f(x)\cdot_T f(r))$. Multiplication preservation for $x\cdot_S r$, followed by target classifier symmetry and ideal transport, gives $J(f(x\cdot_S r))$. Hence $K(x\cdot_S r)$.

The displayed carrier support, zero, additive, inverse, multiplicative, classifier-transport, and two-sided absorption rows are exactly the $\IdealUp$ certificate rows listed in \autoref{sec:ideal-certificate} for a predicate over the source $\RingUp$ certificate.
\end{proof}

\begin{theorem}[Ring map ideal preimage quotient-kernel export surface]
\label{thm:ideal-preimage-quotientkernel-export-surface}
Let $S$ and $T$ carry $\RingUp$ certificates as in \autoref{ch:concrete-instances-ring-namecert}, with notation as in \autoref{def:ring-map-ideal-preimage-predicate}. Let $f:S\to T$ satisfy the carrier-preservation, classifier-preservation, zero-preservation, addition-preservation, additive-inverse-preservation, and multiplication-preservation hypotheses of \autoref{thm:ring-map-ideal-preimage-ideal-closure}. Let $J$ be an $\IdealUp$ predicate over $T$ with the carrier support, zero membership, additive closure, additive inverse closure, internal multiplication closure, classifier transport, and two-sided target absorption rows listed in \autoref{sec:ideal-certificate}. For $K:=\mathsf{Preim}^{\mathrm{ideal}}_{f,J}$, let $Q_K$ be the quotient-kernel source package of \autoref{def:ideal-quotient-kernel-source-package}. Then $Q_K$ exports the five quotient-free rows consumed by $\QuotientRingUp$:
\begin{enumerate}
\item diagonal exactness for every carried source endpoint $x$:
$$
C_S(x)\Longrightarrow Q_K(x,x);
$$
\item endpoint transport along the source classifier:
$$
Q_K(x,y)\land x\sim_S x'\land y\sim_S y'\Longrightarrow Q_K(x',y');
$$
\item additive representative descent:
$$
Q_K(a,a')\land Q_K(b,b')\Longrightarrow Q_K(a+_S b,a'+_S b');
$$
\item negation representative descent:
$$
Q_K(a,a')\Longrightarrow Q_K(-_S a,-_S a');
$$
\item multiplication representative descent:
$$
Q_K(a,a')\land Q_K(b,b')\Longrightarrow Q_K(a\cdot_S b,a'\cdot_S b').
$$
\end{enumerate}
Thus a ring-map ideal preimage can be passed to the quotient-ring construction through the same quotient-kernel interface as any source-side $\IdealUp$ predicate.
\end{theorem}
\begin{proof}
Apply \autoref{thm:ring-map-ideal-preimage-ideal-closure} to the displayed hypotheses. This proves that $K$ is an $\IdealUp$ predicate over the source $\RingUp$ certificate on $S$.

Instantiate \autoref{thm:ideal-quotientring-export-surface} with the source ring $S$ and the ideal predicate $K$. The quotient-kernel source package in that theorem is exactly the relation $Q_K$ obtained from \autoref{def:ideal-quotient-kernel-source-package}.

Reading off the five displayed components gives diagonal exactness, endpoint transport, additive representative descent, negation representative descent, and multiplication representative descent for $Q_K$. No quotient carrier or preimage-specific quotient primitive is introduced; the consumer sees only the quotient-free kernel rows listed above.
\end{proof}

\begin{theorem}[Finite preimage meet quotient-kernel export]
\label{thm:ideal-finite-preimage-meet-quotient-kernel-export}
Let $S$ carry a $\RingUp$ certificate with carrier $C_S$, classifier $\sim_S$, addition $+_S$, additive inverse $-_S$, zero $0_S$, and multiplication $\cdot_S$. Let $xs:\mathsf{HistSpine}$ be finite. For every $i:\mathsf{Pos}(xs)$, let $T_i$ carry a $\RingUp$ certificate with target carrier $C_i$, classifier $\sim_i$, zero $0_i$, addition $+_i$, additive inverse $-_i$, and multiplication $\cdot_i$; let $f_i:S\to T_i$ be a carried ring map in the sense of \autoref{thm:ring-map-ideal-preimage-ideal-closure}; and let $J_i$ be an $\IdealUp$ predicate over $T_i$. Define the source preimage rows
$$
P_i(x)\Longleftrightarrow \mathsf{Preim}^{\mathrm{ideal}}_{f_i,J_i}(x)
$$
and their finite source meet
$$
K(x)\Longleftrightarrow \mathsf{FinIdealMeet}_{C_S,xs,P}(x).
$$
Then $K$ is an $\IdealUp$ predicate over $S$, and the predicate $K$ supplies the quotient-free kernel surface consumed by $\QuotientRingUp$ through \autoref{thm:ideal-quotientring-export-surface}.
\end{theorem}
\begin{proof}
Fix a position $i:\mathsf{Pos}(xs)$. The hypotheses on $f_i$ are exactly the carrier, classifier, zero, addition, inverse, and multiplication preservation rows required by \autoref{thm:ring-map-ideal-preimage-ideal-closure}. Applying that theorem to $f_i$ and $J_i$ gives an $\IdealUp$ certificate for the source predicate $P_i$ over the fixed ring $S$.

The preceding construction is pointwise in $i$, so every row in the finite spine is now an $\IdealUp$ predicate over the same source $\RingUp$ certificate. Apply \autoref{thm:finite-ideal-intersection-closure} to the family $P$ to obtain an $\IdealUp$ certificate for $K$.

Finally instantiate \autoref{thm:ideal-quotientring-export-surface} with this resulting predicate $K$. The export theorem returns the diagonal, endpoint transport, additive descent, negation descent, and multiplication descent rows for the kernel relation generated by $K$, which is the quotient-free surface required before $\QuotientRingUp$ forms its quotient carrier.
\end{proof}

\begin{theorem}[Ideal quotient-kernel absorption closure]
\label{thm:ideal-quotient-kernel-absorption-closure}
Let $R$ carry a $\RingUp$ certificate with carrier $C$, classifier $\sim_R$, addition $+_R$, additive inverse $-_R$, zero $0_R$, and multiplication $\cdot_R$. Let $I$ be an $\IdealUp$ carried predicate over $R$, and define
$$
x-_R y:=x+_R(-_R y),
\qquad
Q_I(x,y)\Longleftrightarrow C(x)\land C(y)\land I(x-_R y).
$$
Assume the ring certificate supplies multiplication carrier closure, carrier transport, classifier symmetry, subtraction congruence, and the two absorption-difference rows
$$
(r\cdot_R x)-_R(r\cdot_R y)\sim_R r\cdot_R(x-_R y)
$$
and
$$
(x\cdot_R r)-_R(y\cdot_R r)\sim_R (x-_R y)\cdot_R r.
$$
Assume the ideal certificate supplies classifier transport and left and right absorption by every carried ring endpoint. Then
$$
C(r)\land Q_I(x,y)
\Longrightarrow
Q_I(r\cdot_R x,r\cdot_R y)\land Q_I(x\cdot_R r,y\cdot_R r).
$$
\end{theorem}
\begin{proof}
Unfold the displayed definition of $Q_I$. From $Q_I(x,y)$ obtain $C(x)$, $C(y)$, and $I(x-_R y)$. The multiplication carrier row applied to $C(r)$ with $C(x)$ and $C(y)$ gives the carried endpoints $C(r\cdot_R x)$, $C(r\cdot_R y)$, $C(x\cdot_R r)$, and $C(y\cdot_R r)$.

For the left-multiplied pair, ideal left absorption gives $I(r\cdot_R(x-_R y))$ from $C(r)$ and $I(x-_R y)$. The left absorption-difference row compares $(r\cdot_R x)-_R(r\cdot_R y)$ with $r\cdot_R(x-_R y)$. Classifier symmetry orients this comparison from the absorbed representative difference to the displayed product difference, and ideal classifier transport yields
$$
I((r\cdot_R x)-_R(r\cdot_R y)).
$$

For the right-multiplied pair, ideal right absorption gives $I((x-_R y)\cdot_R r)$. The right absorption-difference row compares $(x\cdot_R r)-_R(y\cdot_R r)$ with $(x-_R y)\cdot_R r$. Again classifier symmetry orients the comparison for the ideal transport row, giving
$$
I((x\cdot_R r)-_R(y\cdot_R r)).
$$

Pair the first transported membership with $C(r\cdot_R x)$ and $C(r\cdot_R y)$ to obtain $Q_I(r\cdot_R x,r\cdot_R y)$. Pair the second transported membership with $C(x\cdot_R r)$ and $C(y\cdot_R r)$ to obtain $Q_I(x\cdot_R r,y\cdot_R r)$. These two witnesses are exactly the claimed ambient left and right multiplication closure for the ideal quotient-kernel relation.
\end{proof}

\begin{closurestatus}{\IdealUp}
  \theoryclosure{\scopedClosure}
  \scopeclosed{The RingUp-scoped $BHist$ ideal predicate is closed for carrier support, zero, additive, inverse, internal multiplication, classifier transport, two-sided ambient absorption, finite meet, sum, preimage, principal generated ideal, and quotient-kernel export surface.}
  \formalstatus{\theoremCheckedV}
  \leantarget{BEDC.Derived.IdealUp.IdealQuotientKernel\_export\_surface}
  \bridgestatus{none}
  \notclaimed{This chapter does not close quotient carriers, standard ideal theory beyond the displayed predicate surface, or any host quotient primitive.}
  \upgradepath{Public closure requires a public IdealUp interface theorem consumed by QuotientRingUp and IdealClassUp, together with an explicit standard bridge.}
\end{closurestatus}
